% Gauss, Werke X/1 (1917), native106–107, printed99–100.
\SourceLeaf{106}\NoSourceHead
\pdfbookmark[1]{Briefwechsel zum Fundamentalsatz der Algebra}{algebra-briefwechsel}
\HistoricTitle{BRIEFWECHSEL\\[6pt][ZUM FUNDAMENTALSATZ DER ALGEBRA.]}
\NumberedItem{1}
\begin{center}\large\textsc{Pfaff} an \textsc{Gauss}.\end{center}
\begin{flushright}\{Helmstedt 30. May 1799\end{flushright}
\begin{center}Ew. Wohlgeb.\end{center}
\noindent\makebox[\linewidth][s]{. . . . . . . . . . . . . . . . . . . . . . . . . . .}\par
\noindent\hspace*{\parindent}Verschiedene Kleinigkeiten in Ihrer Schrift, die verschrieben oder in der
Eile versehen sind, habe ich mir bereits angemerkt, um sie Ihnen beym Zurücksenden mitzutheilen,
damit sie nicht etwa beym Abdruck zufällig, wenn die Gedanken auf das wichtigere geheftet sind,
wieder übersehen werden. Z. B. § 18 steht: simili modo $T$ ubique \underline{inter 7 et 9},
inter \underline{15 et 17} etc. et generaliter inter $\underline{8k+7}$ et $\underline{8k+9}$
valorem negativum habebit; wo die unterstrichenen Worte heissen müssen: inter 5 et 7,
13 et 15, $8k+5$ et $8k+7$. Auch vorher: $T$ .\,.\,. positivum valorem habebit
.\,.\,.\,.\,.\,. a $8k+3$ usque ad $8k+5$, anstatt $8k+1$ ad $8k+3$.
Einige weitere Anmerkungen über die Sache und Ihre Darstellung wollte ich Ihnen auch schreiben:
es ist mir aber, besonders durch eine noch diesen Nachmittag dazwischen gekommene nicht wohl
zu vermeidende Abhaltung, die Zeit für heute zu kurz geworden, und doch wollte ich Sie nicht
noch einen Posttag länger auf Antwort warten lassen, da Sie ausdrücklich Beschleunigung der
ganzen Angelegenheit gewünscht haben. Eines an sich unbedeutenden Umstandes muss ich jedoch
noch erwähnen. Sie werden selbst zugeben, dass die Betrachtung der krummen Linien und ihrer
Schenkel, auf deren nothwendigem Durchschnitt Ihr Beweis beruht,
\par\vfill\noindent\hfill 13\textsuperscript{*}\par
\clearpage
\SourceLeaf{107}\ArithmeticHead{100}{L}{ALGEBRA.\quad BRIEFWECHSEL.}
\noindent etwas verwickelt und abstract ist. Sollte es daher nicht nützlich und zur Erläuterung
dienlich seyn, wenn Sie etwa für einen Fall in concreto, für welchen der zu beweisende Satz schon
aus anderen Gründen ausgemacht ist, wie $m=4$, die Natur dieser krummen Linien genauer
entwickelten, und durch Zeichnung darstellten, wobey Sie auch noch allenfalls die Coëfficienten
$A,B,C,\,.\,.\,.$ zweckmässig auswählen könnten? Es würde dieses wohl kein Auswuchs seyn,
durch den etwa die Concinnitaet des Ganzen leiden könnte. In den Figuren sind Sie auch fast zu
sparsam gewesen, da Sie die Fig. 2 in verschiedenen Bedeutungen und Absichten gebrauchen,
indem die mit Zahlen bezeichneten Puncte § 18 etwas ganz anders vorstellen, als § 20.
Ich bin der Meynung, dass man beym Schreiben sich selbst die Pflicht aufzulegen hat, auch für
die Bequemlichkeit der Leser zu sorgen, selbst wenn es gewissermassen mit eigener Unbequemlichkeit
verknüpft wäre. Nächst dem materiellen Gehalt sind doch immer Ordnung und Deutlichkeit die
ersten formellen Eigenschaften jedes Vortrags. Gegen die Ordnung, worin Sie Ihre Ideen
entwickeln, habe ich nichts einzuwenden, aber der Deutlichkeit möchte für die Mehrheit der Leser
noch etwas nachzuhelfen seyn. In beyden ist \textsc{Euler} mehr zum Muster zu empfehlen, als die
neuen französischen Mathematiker. Sie werden diese Anmerkungen für keinen Tadel halten:
ich wünsche aber auch, dass Sie solche nicht für einen blossen Rath, um zu rathen, ansehen,
sondern dadurch für die Zukunft auf einen Punct aufmerksamer werden mögen, der wirklich
zum Vortheil der Wissenschaft wichtiger ist, als man leicht zu glauben geneigt ist.

Ich empfehle mich Ihrem ferneren freundschaftl. Andenken, und verharre mit vollkommenster
\begin{flushright}
\begin{tabular}{c}Hochachtung\\Ihr ergebenster Freund\\J. F. \textsc{Pfaff}.\end{tabular}
\end{flushright}
\noindent\hspace*{\parindent}N.S. An H. H[of]r[ath] v. Z[\textsc{immermann}] viele Empfehlungen.
Sie sprachen bey Ihrem Hierseyn von \textsc{Bachet}’s \textit{Problemes plaisans.}
Ich besann mich damals nicht gleich, dass ich das Buch selbst habe, das Ihnen also auf
Verlangen zu Dienst steht.\}
\ClosingRule
\clearpage
